\noindent\textbf{Resumen.} Esta parte explica cómo se reparte el trabajo sobre los 5,7 millones de contratos en tres niveles. Step Functions distribuye 284 tramos entre varias Lambdas, cada Lambda descarga con hilos y Spark, dentro de AWS Glue, limpia y agrega los datos. Con cuatro hilos la descarga completa pasó de 44,5 a 24 minutos ($S = 1{,}86$), y con ocho ya no bajó más. En Spark, subir de 2 a 4 \textit{workers} redujo el tiempo de 277 a 193~s ($S = 1{,}44$).

% ---------------------------------------------------------------------------
\subsection{Introducción}
Bajar y transformar el histórico con un solo proceso toma horas, y los indicadores que comparan cada contrato con los demás de su entidad no caben en la memoria de un portátil. La Sección~\ref{sec:intro} presenta los objetivos y el ODS~16.

Hay dos cuellos de botella. La descarga depende de la red, de los límites de la API y de los 15 minutos que como máximo dura una Lambda. Para los indicadores hay que reunir los contratos de cada entidad, de cada proveedor y de cada clase de contrato, y eso obliga a Spark a mover datos entre nodos (\textit{shuffle}).

Este diseño se apoya en trabajos conocidos. MapReduce mostró cómo procesar lotes en clústeres de máquinas comunes~\cite{dean2008}, y Spark agregó el trabajo en memoria y la reconstrucción de una partición perdida a partir de su linaje~\cite{zaharia2012,zaharia2016}. El \textit{lakehouse} lleva transacciones ACID al almacenamiento de objetos~\cite{armbrust2021}. La idea de medir riesgo en compras públicas con indicadores objetivos viene de Fazekas y Kocsis~\cite{fazekas2020}.

% ---------------------------------------------------------------------------
\subsection{Sistema propuesto}
En la Fig.~\ref{cpd:fig:arq} aparece el recorrido completo. AWS administra todos los servicios, así que el equipo no mantiene servidores y paga solo el tiempo de cómputo. Bronce guarda lo que responde la API, plata los datos limpios y oro las seis tablas que se consultan con Athena.

\begin{figure}[!htbp]
\centering



\begin{tikzpicture}[
  node distance=4mm and 3mm,
  box/.style={draw, rounded corners, align=center, font=\scriptsize, minimum height=7mm, text width=17mm},
  layer/.style={draw, fill=#1, rounded corners, align=center, font=\scriptsize, minimum height=7mm, text width=17mm},
  arr/.style={-{Stealth[length=2mm]}, thick}
]
\node[box] (api) {Fuente\\SECOP~II};
\node[box, right=of api] (lam) {Descarga\\12 en paralelo};
\node[layer=orange!25, right=of lam] (bro) {Bronce\\(crudo)};
\node[box, below=of bro] (g1) {Spark:\\limpieza};
\node[layer=gray!20, left=of g1] (pla) {Plata\\(limpio)};
\node[box, left=of pla] (g2) {Spark:\\indicadores};
\node[layer=yellow!35, below=of g2] (oro) {Oro\\(resúmenes)};
\node[box, right=of oro] (ath) {Consultas\\SQL};
\draw[arr] (api) -- (lam);
\draw[arr] (lam) -- (bro);
\draw[arr] (bro) -- (g1);
\draw[arr] (g1) -- (pla);
\draw[arr] (pla) -- (g2);
\draw[arr] (g2) -- (oro);
\draw[arr] (oro) -- (ath);
\node[draw, dashed, rounded corners, fit=(api)(lam)(bro)(g1)(pla)(g2), inner sep=2mm,
  label={[font=\scriptsize]above:Coordinado cada semana por Step Functions}] {};
\end{tikzpicture}
\caption{Recorrido de los datos por las capas bronce, plata y oro.}
\label{cpd:fig:arq}
\end{figure}

% ---------------------------------------------------------------------------
\subsection{Diseño paralelo y distribuido}

\subsubsection{Nivel 1: tramos en Lambdas concurrentes}
Step Functions, con un estado \texttt{Map}, entrega los tramos de tiempo a las Lambdas de tres en tres. Ningún tramo necesita datos de otro, de modo que el problema es \emph{embarrassingly parallel}.

El tamaño del tramo importa. Enero de 2026 tiene 487.368 contratos y, con un hilo, habría tardado unos 1.740~s, casi el doble de los 900~s que permite una Lambda. Por eso los cinco meses más grandes se bajan por días, y el total queda en 284 tareas: 133 meses, 150 días y una para los contratos sin fecha. Se dejaron tres Lambdas a la vez porque la API limita las peticiones por token, y con más solo subían los errores.

\subsubsection{Nivel 2: hilos dentro de cada Lambda}
Los hilos de un mismo proceso usan la misma memoria y los mismos recursos, como el cliente de S3, y crearlos cuesta menos que crear procesos~\cite{grama2003}. Casi todo el tiempo de una página se va en esperar la respuesta de la API, y con varios hilos esas esperas ocurren al mismo tiempo (Fig.~\ref{cpd:fig:hilos}).

\begin{figure}[!htbp]
\centering



\begin{tikzpicture}[x=6.2mm, y=4.2mm, font=\scriptsize]
  \newcommand{\pag}[3]{%
    \fill[gray!35] (#1,#2) rectangle ++(1.6,0.7);
    \fill[black!80] (#1+1.6,#2) rectangle ++(0.2,0.7);
    \fill[orange!60] (#1+1.8,#2) rectangle ++(0.4,0.7);
    \node[font=\tiny] at (#1+0.8,#2+0.35) {#3};}
  \node[anchor=east] at (-0.1,7.35) {1 hilo};
  \foreach \i in {0,...,4} { \pag{\i*2.2}{7}{p\i} }
  \node[anchor=east] at (-0.1,5.35) {Hilo 1};
  \node[anchor=east] at (-0.1,4.35) {Hilo 2};
  \node[anchor=east] at (-0.1,3.35) {Hilo 3};
  \node[anchor=east] at (-0.1,2.35) {Hilo 4};
  \pag{0}{5}{p0} \pag{2.2}{5}{p4}
  \pag{0}{4}{p1} \pag{0}{3}{p2} \pag{0}{2}{p3}
  \draw[-{Stealth[length=1.5mm]}] (0,1.4) -- (11.4,1.4) node[right] {$t$};
  \draw[dashed] (11,1.3) -- (11,7.9);
  \draw[dashed] (4.4,1.3) -- (4.4,6.0);
  \node[below] at (4.4,1.4) {$T_4$};
  \node[below] at (11,1.4) {$T_1$};
  \fill[gray!35] (0,0) rectangle ++(0.5,0.5); \node[right] at (0.5,0.25) {espera de la API};
  \fill[black!80] (4.6,0) rectangle ++(0.5,0.5); \node[right] at (5.1,0.25) {CPU};
  \fill[orange!60] (6.9,0) rectangle ++(0.5,0.5); \node[right] at (7.4,0.25) {escritura en S3};
\end{tikzpicture}
\caption{Cinco páginas bajadas con uno y con cuatro hilos. Con cuatro, los tiempos de espera coinciden y la quinta página sale cuando el primer hilo se desocupa.}
\label{cpd:fig:hilos}
\end{figure}

En CPython un solo hilo ejecuta código Python en cada instante, por el GIL~\cite{psf_threading,psf_futures}. El GIL se suelta mientras el hilo espera la red o comprime, que es lo que más tarda aquí. Lo que sí se turna, como convertir el JSON, hace de fracción secuencial $f$ en la ley de Amdahl. Si el trabajo fuera sobre todo cálculo, convendría usar procesos.

Para repartir, la Lambda cuenta los registros del tramo, calcula dónde empieza cada página y pasa esa lista a un \texttt{ThreadPoolExecutor}:

\begin{lstlisting}[language=Python]
total = count_rows(dataset_id, where)
offsets = range(0, total, page_size)
with ThreadPoolExecutor(max_workers=workers) as pool:
    sizes = list(pool.map(job, enumerate(offsets)))
\end{lstlisting}

Cada hilo recibe su propio \textit{offset} y escribe un archivo con otro nombre, por lo que no hacen falta candados. El programa no sale del bloque \texttt{with} hasta que terminan todos los hilos, y solo entonces el hilo principal suma los conteos. Con cuatro hilos la memoria subió a 2.048~MB, porque cada uno guarda su página mientras la sube a S3. La variable \texttt{PARALLEL\_\allowbreak PAGES} fija cuántos hilos usar; con uno, el programa equivale a la versión secuencial.

El tiempo de un tramo con $P$ páginas que tardan $t$ cada una, un costo fijo $c$ y $p$ hilos es aproximadamente
\begin{equation}
T_p \approx c + \left\lceil \frac{P}{p} \right\rceil t, \qquad S(p) = \frac{T_1}{T_p} \leq \frac{P}{\lceil P/p \rceil}.
\end{equation}
El límite lo pone el reparto de páginas. Con 9 páginas, cuatro hilos llegan como mucho a $S = 3$ y ocho a $S = 4{,}5$, y un tramo de una sola página tarda lo mismo con cualquier número de hilos.

Entre Lambdas no hay memoria compartida, así que ese paralelismo es distribuido. Dentro de cada Lambda los hilos comparten memoria. En total puede haber $3 \times 4 = 12$ descargas al mismo tiempo.

\subsubsection{Nivel 3: clúster de Spark}
El trabajo \texttt{bronze\_to\_silver} lee todos los archivos a la vez, limpia, deja una fila por \texttt{id\_contrato} con \texttt{row\_number()} y escribe con \texttt{MERGE INTO}. El trabajo \texttt{silver\_to\_gold} calcula las nueve alertas de la Tabla~\ref{adq:t7}. Varias de ellas obligan a un \textit{shuffle}: los percentiles por tipo de contrato para el valor atípico ($Q_3 + 3\,\mathrm{IQR}$), las ventanas de 30 días por entidad, contratista y código UNSPSC para el fraccionamiento, y la suma por entidad para la concentración. El límite del anticipo sale de la Ley 80~\cite{ley80}.

\subsubsection{Tolerancia a fallos y escalabilidad}
Cada petición tiene cuatro intentos y cada tramo, tres; durante la carga histórica fallaron cuatro tramos y todos se recuperaron sin ayuda. Si un nodo de Spark cae, se recalculan únicamente sus particiones~\cite{zaharia2012}. El \texttt{MERGE} no crea duplicados e Iceberg confirma cada escritura completa o no la confirma~\cite{asf_iceberg_spec}; dos trabajos que se detuvieron a propósito dejaron la tabla igual. Para crecer basta con subir los tramos simultáneos, los hilos o los \textit{workers}, que tienen 4~vCPU y 16~GB cada uno, y plata está dividida por mes para leer solo lo necesario.

% ---------------------------------------------------------------------------
\subsection{Metodología de medición}
Si $T_1$ es el tiempo con un hilo o proceso y $T_p$ el tiempo con $p$~\cite{grama2003}, entonces
\begin{equation}
S(p) = \frac{T_1}{T_p}, \qquad E(p) = \frac{S(p)}{p}, \qquad S(p) \leq \frac{1}{f + (1-f)/p}.
\end{equation}
La tercera expresión es la ley de Amdahl~\cite{amdahl1967}: si una fracción $f$ es secuencial, el \textit{speedup} tiene techo. Según la ley de Gustafson~\cite{gustafson1988}, ese techo sube cuando el tamaño de los datos crece junto con los recursos.

En el experimento A se repitió la descarga completa (5.698.182 contratos en 371 tramos) con 1, 2, 4 y 8 hilos, siempre con tres Lambdas. Las corridas se hicieron una tras otra, hacia una carpeta aparte, y el tiempo se tomó de Step Functions. En el experimento B se corrió \texttt{silver\_to\_gold} con 2, 4, 8 y 10 \textit{workers} en modo Standard, en tres series con distinto orden, y se reporta la mediana. Como Glue usa un \textit{worker} como \textit{driver}, la base es $T_2$ y la eficiencia es $E(n) = S(n)/(n-1)$. Los scripts \texttt{benchmark\_historico.py} y \texttt{benchmark\_spark.py} permiten repetir ambos experimentos.

% ---------------------------------------------------------------------------
\subsection{Resultados}

\subsubsection{Carga histórica}
Los primeros 74 meses se bajaron con un hilo y el resto con cuatro (Tabla~\ref{cpd:tab:fase}). Al terminar, bronce tenía 5.687.382 contratos en 732 archivos.

\begin{table}[!htbp]
\caption{Segunda parte de la carga histórica, ya con hilos}
\label{cpd:tab:fase}
\centering
\begin{tabular}{@{}lr@{}}
\toprule
Tramos completados & 210 de 210 \\
Contratos descargados & 4.672.446 \\
Duración real & 27~min~11~s \\
Tiempo sumado de todas las Lambdas & 4.592~s \\
Fallos recuperados / tramos perdidos & 4 / 0 \\
Mediana, tramos de 4 o más páginas & 2.094 reg./s \\
Mediana, tramos de 1 página & 534 reg./s \\
\bottomrule
\end{tabular}
\end{table}

Dividir el tiempo sumado de las Lambdas entre el tiempo de pared da $4.592 / 1.631 \approx 2{,}8$, cerca del máximo de tres Lambdas ocupadas. Los tramos de una sola página avanzaron cuatro veces más lento que los grandes, como anticipa la Ec.~(1). El tramo más pesado, enero de 2022 con 294.128 contratos, tardó 250~s, muy por debajo de los 900~s permitidos.

\subsubsection{Spark: plata y oro}
Plata leyó 5.690.495 filas y conservó 5.687.110 contratos únicos, y las reglas críticas se cumplieron (Tabla~\ref{adq:t6}). Frente a la prueba de un mes, el histórico tenía 68 veces más datos y se procesó en un tiempo parecido con cinco veces más \textit{workers} (Tabla~\ref{cpd:tab:debil}). Pasa porque el arranque del clúster cuesta casi lo mismo con pocos o muchos datos; es escalabilidad débil, en el sentido de Gustafson~\cite{gustafson1988}.

\begin{table}[!htbp]
\caption{Spark: un mes frente al histórico completo}
\label{cpd:tab:debil}
\centering
\begin{tabular}{@{}lrr@{}}
\toprule
 & Sept. 2026 & Histórico \\
\midrule
Contratos & 83.105 & 5.687.110 \\
\textit{Workers} G.1X & 2 & 10 \\
Tiempo (s) & 208 & 199 \\
Contratos por segundo y \textit{worker} & $\approx$ 200 & $\approx$ 2.860 \\
\bottomrule
\end{tabular}
\end{table}

Oro consume menos que plata con la mitad de \textit{workers}, porque lee Parquet con pocas columnas y plata lee JSON con 91. En modo Flex, dos corridas de plata sobre los mismos datos tardaron 188 y 240~s, mientras que sus DPU-s apenas variaron un 5~\%. La diferencia estuvo en la espera por \textit{workers}, y por eso el experimento~B se corrió en modo Standard. Las alertas y la Ley de Garantías están en las Tablas~\ref{adq:t7} y \ref{adq:t8}.

\subsubsection{Carga semanal}
Cada lunes a las 6:00~a.~m., EventBridge Scheduler lanza el pipeline. Como la API publica con días de retraso, se vuelven a bajar los últimos 17 días y el \texttt{MERGE} actualiza lo que ya existía. Hasta el 23 de noviembre de 2026 la carga era diaria y bajaba 10 días (Tabla~\ref{cpd:tab:diario}); cada ejecución cuesta unos 0,07~USD.

\begin{table}[!htbp]
\caption{Una actualización completa, medida el 13 nov. 2026, cuando la carga era diaria}
\label{cpd:tab:diario}
\centering
\begin{tabular}{@{}lrr@{}}
\toprule
Paso & \textit{Workers} / hilos & Tiempo (s) \\
\midrule
Lambda (25.533 contratos) & 4 hilos & 22 \\
Bronce $\rightarrow$ plata & 2 & 129 \\
Plata $\rightarrow$ oro & 5 & 230 \\
\midrule
Ejecución completa & & 491 \\
\bottomrule
\end{tabular}
\end{table}

\subsubsection{Experimento A: hilos}
Las cuatro corridas terminaron sin errores y produjeron los mismos 784 archivos (Tabla~\ref{cpd:tab:mp}). La columna $e$ aplica la métrica de Karp y Flatt~\cite{karp1990}, que deduce qué parte del trabajo sigue siendo secuencial a partir de la aceleración observada:
\begin{equation}
e = \frac{1/S(p) - 1/p}{1 - 1/p}.
\end{equation}

\begin{table}[!htbp]
\caption{Experimento A: el histórico completo con distinto número de hilos}
\label{cpd:tab:mp}
\centering
\begin{tabular}{@{}rrrrrr@{}}
\toprule
Hilos $p$ & $T_p$ (s) & $S(p)$ & $E(p)$ & $\Sigma$ Lambdas (s) & $e$ \\
\midrule
1 & 2.673 & 1,00 & 1,00 & 7.902 & -- \\
2 & 1.565 & 1,71 & 0,85 & 4.598 & 0,17 \\
4 & 1.439 & 1,86 & 0,47 & 4.234 & 0,38 \\
8 & 1.492 & 1,79 & 0,22 & 4.378 & 0,50 \\
\bottomrule
\end{tabular}
\end{table}

Si la parte secuencial fuera fija, $e$ no cambiaría; aquí pasa de 0,17 a 0,50, de modo que cada hilo extra agrega su propio costo. La Tabla~\ref{cpd:tab:granularidad} separa los tramos por tamaño para ver de dónde sale.

\begin{table}[!htbp]
\caption{Segundos de Lambda gastados según el tamaño del tramo}
\label{cpd:tab:granularidad}
\centering
\begin{tabular}{@{}lrrrr@{}}
\toprule
Tramos & $p=1$ & $p=2$ & $p=4$ & $p=8$ \\
\midrule
252 de 1 página & 1.435 & 1.420 & 1.498 & 1.670 \\
53 de 5 o más páginas & 4.581 & 2.149 & 1.655 & 1.585 \\
\bottomrule
\end{tabular}
\end{table}

El 68~\% de los tramos tiene una sola página, así que no gana nada con más hilos, y con ocho tarda un 16~\% más por crear hilos sin trabajo. Los 53 tramos grandes sí aprovechan los hilos y llegan a $S = 2{,}9$ con ocho. Con tres Lambdas de ocho hilos hay hasta 24 peticiones a la vez contra la misma API, que pasa a ser el límite. Con cuatro hilos se obtuvo el mejor \textit{speedup} y el menor tiempo sumado de Lambda, el 54~\% del de un hilo, que es también el menor costo. Esa es la configuración que quedó en producción.

\subsubsection{Experimento B: Spark}
Se completaron doce corridas sobre 5.708.737 contratos (Tabla~\ref{cpd:tab:spark}).

\begin{table}[!htbp]
\caption{Experimento B: cálculo de oro con distinto número de \textit{workers}}
\label{cpd:tab:spark}
\centering
\begin{tabular}{@{}rrrrrr@{}}
\toprule
\textit{Workers} $n$ & $T_n$ (s) & Rango (s) & $S(n)$ & $E(n)$ & $e$ \\
\midrule
2 & 277 & 255 a 293 & 1,00 & 1,00 & -- \\
4 & 193 & 183 a 194 & 1,44 & 0,48 & 0,55 \\
8 & 171 & 163 a 207 & 1,62 & 0,23 & 0,55 \\
10 & 213 & 194 a 222 & 1,30 & 0,14 & 0,74 \\
\bottomrule
\end{tabular}
\end{table}

Entre 2 y 4 \textit{workers} los rangos de tiempo no se tocan, así que la mejora es real. De 4 en adelante las diferencias caben dentro de lo que varía de una serie a otra. Con 10 \textit{workers} $e$ sube a 0,74, lo que apunta a un costo de coordinación en el arranque, el \textit{shuffle} o la escritura de las seis tablas; el experimento no permite separar esas causas. Con 4 \textit{workers} el tiempo baja un 30~\% a cambio de un 40~\% más de DPU-s, y por encima de eso se paga más sin ganar tiempo.

% ---------------------------------------------------------------------------
\subsection{Costos}
Un presupuesto de 15~USD avisa al 50, 80 y 100~\% y, al llegar al tope, bloquea Lambda, Step Functions y Glue. Glue en modo Flex costó unos 0,83~USD en todo el proyecto. El experimento~B costó 1,78~USD y el~A, cerca de 1~USD.

% ---------------------------------------------------------------------------
\subsection{Conclusiones}
\begin{itemize}
  \item Las Lambdas, los hilos y Spark trabajan juntos como tres niveles de paralelismo sobre el mismo flujo.
  \item Con cuatro hilos la descarga bajó de 44,5 a 24 minutos. Más hilos no ayudan, porque la mayoría de los tramos tiene una página y la API es compartida.
  \item Para este volumen, 4 \textit{workers} de Spark dan el mejor balance entre tiempo y costo.
  \item Los cuatro fallos de la carga histórica se recuperaron sin perder datos, y cada carga semanal cuesta unos 0,07~USD.
  \item Falta medir cada etapa de Spark por separado para saber qué frena su \textit{speedup}.
\end{itemize}
